% TOP_PROBLEM: {"id": 3229, "title": "Coloring and immersion", "category_id": 3, "category": {"id": 3, "name": "graph_theory", "display_name": "Graph Theory", "description": "Problems involving graphs, networks, and their properties.", "slug": "graph-theory", "order_index": 3, "created_at": "2026-07-31T15:26:25.671Z"}, "status": "open", "problem_number": "OPG-550", "set_id": 12, "statusQualification": "Specifies weak immersion by pairwise edge-disjoint paths, avoiding a typographical error in the source’s edge-splitting description."}
% FIRST_POSTED: 2026-10-01
% ENTRY_KIND: statement_only
% UnsolvedMath ID 3229; source catalogue code OPG-550
% !TeX program = lualatex
% Definition and sources only; this file is not an LLM solution attempt.
\documentclass[11pt,a4paper]{article}
\usepackage[margin=25mm]{geometry}
\usepackage{fontspec}
\setmainfont{lmroman10-regular.otf}[BoldFont=lmroman10-bold.otf,
 ItalicFont=lmroman10-italic.otf,BoldItalicFont=lmroman10-bolditalic.otf]
\usepackage{amsmath,amssymb,mathtools}
\usepackage{unicode-math}
\setmathfont{latinmodern-math.otf}
\AtBeginDocument{\let\setminus\smallsetminus}
\usepackage{xurl,hyperref}
\hypersetup{colorlinks=true,urlcolor=blue}
\setcounter{secnumdepth}{0}
\setlength{\parindent}{0pt}
\setlength{\parskip}{5pt}
\setlength{\emergencystretch}{3em}
\begin{document}
\section*{Coloring and immersion}\label{3229}
{\small UnsolvedMath ID 3229 · OPG-550 · Graph Theory · OpenGarden}
\subsection{Definitions and mathematical statement}
Let $G$ be a finite loopless undirected graph; parallel edges may be retained as distinct edges. A proper vertex coloring assigns different colors to adjacent vertices, and $\chi(G)$ is the minimum number of colors required. For an integer $t\ge1$, an immersion of $K_t$ in $G$ consists of distinct vertices $b_1,\ldots,b_t$ and, for each unordered pair $\{i,j\}$, a path $P_{ij}$ from $b_i$ to $b_j$ such that all these paths are pairwise edge-disjoint.

This uses the weak immersion convention: different paths may meet at internal vertices, and a branch vertex $b_i$ may occur internally on another pair's path. A strong immersion would additionally exclude that latter possibility and is a different assertion. The paths may otherwise be taken simple, since cycles can be removed from an individual path.

The conjecture states that every finite graph $G$ contains an immersion of $K_{\chi(G)}$. Equivalently, for every positive integer $t$, a graph with chromatic number at least $t$ should have an immersed complete graph of order $t$. Extra edges and vertices of $G$ can be ignored.

The edge-disjointness requirement distinguishes immersion from a complete minor, whose connected branch sets must be vertex-disjoint. Immersion also differs from a subdivision, which requires the internal vertices of distinct edge paths to be disjoint. The target here is the edge-disjoint path model above, with its branch-vertex convention fixed explicitly so that these related notions cannot be confused.
\subsection{Short English statement}
Must every graph contain a weak immersion of a complete graph whose order equals its chromatic number?
\subsection{Sources}
\begin{itemize}
\item[S1] ULAM.AI and the UnsolvedMath Contributors, supplied problems.json, record 3229 (OPG-550), OpenGarden. Catalogue identity and source transcription; CC BY 4.0.
\url{https://huggingface.co/datasets/ulamai/UnsolvedMath}.
\item[S2] Open Problem Garden, Coloring and immersion. Original problem and discussion.
\url{http://www.openproblemgarden.org/op/coloring_and_immersion}.
\end{itemize}
\end{document}
